% ============================================================================
% payments-bench -- project overview, rendered to PDF.
%
% This is NOT a research paper under architecture-standards'
% docs/research/00-RESEARCH-DOCUMENTATION.md. That standard governs studies
% that answer a previously unknown question with evidence; this document is a
% project overview for a reader who wants the whole shape of the repository in
% one sitting. It \inputs the house preamble (house-preamble.tex, copied
% unmodified from architecture-standards) purely for the house look, and
% repurposes \paperstatus: it carries the document's kind and the state of the
% project, not a draft/verified research marker.
%
% Source of truth: docs/OVERVIEW.md in this repository. This file is a curated
% presentation of that document and introduces no fact it does not contain.
% Nothing enforces that an edit to OVERVIEW.md reaches this file; that drift
% is an accepted consequence (docs/plans/stream-c-docs.md), to be met with a
% section-heading check if the two are ever found to disagree.
%
% The partner bank in every section is SIMULATED
% (src/PaymentsBench.PartnerSimulator). Nothing here has talked to a real bank.
%
% Figures are the repository's own Mermaid sources (docs/diagrams/*.mmd),
% rendered to vector PDF by `npm run render:diagrams` before this file is
% compiled. ASCII-only source, so pdflatex is enough.
%
% Build (PDFs are build output -- never committed):
%   npm ci && npm run render:diagrams
%   cd docs/papers && latexmk -pdf payments-bench-overview.tex
% Also built by .github/workflows/build-overview-pdf.yml: on demand, on pull
% requests touching docs/papers/** or docs/diagrams/** (run artifact), and on
% v* tags (attached to the GitHub Release as PaymentsBench_Overview.pdf).
% ============================================================================
\documentclass[11pt,a4paper]{article}

\newcommand{\paperstatus}{PROJECT OVERVIEW}
\newcommand{\headerleft}{payments-bench --- Project overview}
\newcommand{\pdftitleline}{payments-bench: an outbound transfer service that moves money exactly once}
\newcommand{\pdfauthorline}{payments-bench contributors}
% ============================================================================
% HOUSE-PREAMBLE.tex -- the estate's house LaTeX style, as a file you copy in
% and \input rather than a block you copy out.
%
% See 00-RESEARCH-DOCUMENTATION.md, "The house preamble is a file, not a
% block". Copy this file to <repo>/docs/papers/house-preamble.tex (or
% docs/research/papers/ for a study paper) and \input it from every document
% in the repository that wants the house look. One copy per repository; every
% document in that repository \inputs the same one.
%
% WHY A FILE. The estate's first attempt at this shipped a template whose
% preamble was meant to be copied into each new document. Copying is what
% drifts: `ab-ove`'s two overview editions carried none of the seven house
% markers -- not the colors, the section formatting, the header, the status
% marker, the path macro, the geometry or the link style -- and nothing said
% so, because a convention about copying correctly has nothing to fail. The
% Beamer theme, which was always distributed as a file, propagated
% byte-identically over the same period. Distribute the file.
%
% ----------------------------------------------------------------------------
% THE CONTRACT. The including document defines these BEFORE \input, and
% carries its own \documentclass and \begin{document} -- this file has
% neither, so that one preamble can serve documents with different classes.
%
%   \paperstatus    Right-hand header, and the title block's status line.
%                   A study paper: DRAFT or VERIFIED, mirroring the study's
%                   own header. Any other document: its own status, named in
%                   that document's header comment as repurposed.
%   \headerleft     Left-hand header. Usually "Project --- Document title".
%   \pdftitleline   PDF metadata title.
%   \pdfauthorline  PDF metadata author.
%
% Optional, both defaulted below if the document does not define them:
%
%   \houselang      babel language for this edition, e.g. \def\houselang{polish}.
%                   Defaults to english. Set it in the edition file, not here:
%                   it is the one thing that is genuinely per-edition.
%   \editionsuffix  This edition's diagram-file suffix, e.g.
%                   \def\editionsuffix{.pl}. Defaults to empty. \dgm reads it,
%                   so a figure is called by slug and resolves per edition --
%                   see "Figures" below.
%
% A minimal English document:
%
%   \documentclass[11pt,a4paper]{article}
%   \newcommand{\paperstatus}{DRAFT}
%   \newcommand{\headerleft}{Project --- Overview}
%   \newcommand{\pdftitleline}{Project: an overview}
%   \newcommand{\pdfauthorline}{Author Name}
%   % ============================================================================
% HOUSE-PREAMBLE.tex -- the estate's house LaTeX style, as a file you copy in
% and \input rather than a block you copy out.
%
% See 00-RESEARCH-DOCUMENTATION.md, "The house preamble is a file, not a
% block". Copy this file to <repo>/docs/papers/house-preamble.tex (or
% docs/research/papers/ for a study paper) and \input it from every document
% in the repository that wants the house look. One copy per repository; every
% document in that repository \inputs the same one.
%
% WHY A FILE. The estate's first attempt at this shipped a template whose
% preamble was meant to be copied into each new document. Copying is what
% drifts: `ab-ove`'s two overview editions carried none of the seven house
% markers -- not the colors, the section formatting, the header, the status
% marker, the path macro, the geometry or the link style -- and nothing said
% so, because a convention about copying correctly has nothing to fail. The
% Beamer theme, which was always distributed as a file, propagated
% byte-identically over the same period. Distribute the file.
%
% ----------------------------------------------------------------------------
% THE CONTRACT. The including document defines these BEFORE \input, and
% carries its own \documentclass and \begin{document} -- this file has
% neither, so that one preamble can serve documents with different classes.
%
%   \paperstatus    Right-hand header, and the title block's status line.
%                   A study paper: DRAFT or VERIFIED, mirroring the study's
%                   own header. Any other document: its own status, named in
%                   that document's header comment as repurposed.
%   \headerleft     Left-hand header. Usually "Project --- Document title".
%   \pdftitleline   PDF metadata title.
%   \pdfauthorline  PDF metadata author.
%
% Optional, both defaulted below if the document does not define them:
%
%   \houselang      babel language for this edition, e.g. \def\houselang{polish}.
%                   Defaults to english. Set it in the edition file, not here:
%                   it is the one thing that is genuinely per-edition.
%   \editionsuffix  This edition's diagram-file suffix, e.g.
%                   \def\editionsuffix{.pl}. Defaults to empty. \dgm reads it,
%                   so a figure is called by slug and resolves per edition --
%                   see "Figures" below.
%
% A minimal English document:
%
%   \documentclass[11pt,a4paper]{article}
%   \newcommand{\paperstatus}{DRAFT}
%   \newcommand{\headerleft}{Project --- Overview}
%   \newcommand{\pdftitleline}{Project: an overview}
%   \newcommand{\pdfauthorline}{Author Name}
%   % ============================================================================
% HOUSE-PREAMBLE.tex -- the estate's house LaTeX style, as a file you copy in
% and \input rather than a block you copy out.
%
% See 00-RESEARCH-DOCUMENTATION.md, "The house preamble is a file, not a
% block". Copy this file to <repo>/docs/papers/house-preamble.tex (or
% docs/research/papers/ for a study paper) and \input it from every document
% in the repository that wants the house look. One copy per repository; every
% document in that repository \inputs the same one.
%
% WHY A FILE. The estate's first attempt at this shipped a template whose
% preamble was meant to be copied into each new document. Copying is what
% drifts: `ab-ove`'s two overview editions carried none of the seven house
% markers -- not the colors, the section formatting, the header, the status
% marker, the path macro, the geometry or the link style -- and nothing said
% so, because a convention about copying correctly has nothing to fail. The
% Beamer theme, which was always distributed as a file, propagated
% byte-identically over the same period. Distribute the file.
%
% ----------------------------------------------------------------------------
% THE CONTRACT. The including document defines these BEFORE \input, and
% carries its own \documentclass and \begin{document} -- this file has
% neither, so that one preamble can serve documents with different classes.
%
%   \paperstatus    Right-hand header, and the title block's status line.
%                   A study paper: DRAFT or VERIFIED, mirroring the study's
%                   own header. Any other document: its own status, named in
%                   that document's header comment as repurposed.
%   \headerleft     Left-hand header. Usually "Project --- Document title".
%   \pdftitleline   PDF metadata title.
%   \pdfauthorline  PDF metadata author.
%
% Optional, both defaulted below if the document does not define them:
%
%   \houselang      babel language for this edition, e.g. \def\houselang{polish}.
%                   Defaults to english. Set it in the edition file, not here:
%                   it is the one thing that is genuinely per-edition.
%   \editionsuffix  This edition's diagram-file suffix, e.g.
%                   \def\editionsuffix{.pl}. Defaults to empty. \dgm reads it,
%                   so a figure is called by slug and resolves per edition --
%                   see "Figures" below.
%
% A minimal English document:
%
%   \documentclass[11pt,a4paper]{article}
%   \newcommand{\paperstatus}{DRAFT}
%   \newcommand{\headerleft}{Project --- Overview}
%   \newcommand{\pdftitleline}{Project: an overview}
%   \newcommand{\pdfauthorline}{Author Name}
%   \input{house-preamble}
%   \begin{document}
%   ...
%
% Its Polish sibling differs by two lines, both before the \input:
%
%   \def\houselang{polish}
%   \def\editionsuffix{.pl}
%
% ----------------------------------------------------------------------------
% PACKAGES THIS FILE DOES NOT LOAD. The house set below is what every document
% needs. A document that needs more -- tikz, pifont, longtable, fancyvrb,
% amssymb -- loads it in its own preamble AFTER the \input, and says in a
% comment that it is an addition rather than part of the house set, so the
% next reader can tell the two apart at a glance.
% ============================================================================

% ----------------------------------------------------------------------------
% Encoding, then language. fontenc before inputenc before babel: babel reads
% the active font encoding when it sets up its shorthands, so loading it first
% gives an edition that is translated but typeset wrong.
% ----------------------------------------------------------------------------
\usepackage[T1]{fontenc}
\usepackage[utf8]{inputenc}

% The edition's language, defaulted so a monolingual document defines nothing.
% \PassOptionsToPackage rather than \usepackage[\houselang]{babel}: the option
% list here is a macro, and PassOptionsToPackage is the mechanism built for
% passing one.
\providecommand{\houselang}{english}
\PassOptionsToPackage{\houselang}{babel}
\usepackage{babel}

% ----------------------------------------------------------------------------
% The house package set.
% ----------------------------------------------------------------------------
\usepackage{geometry}
\usepackage{xcolor}
\usepackage{hyperref}
\usepackage{graphicx}
\usepackage{enumitem}
\usepackage{fancyhdr}
\usepackage{titlesec}
\usepackage{parskip}
\usepackage{booktabs}
\usepackage{array}
\usepackage{tabularx}

% ----------------------------------------------------------------------------
% Page geometry.
% ----------------------------------------------------------------------------
\geometry{
    a4paper,
    left=2.2cm,
    right=2.2cm,
    top=2.5cm,
    bottom=2.5cm
}

% ----------------------------------------------------------------------------
% House colors. Extracted from the estate's first formal documents
% (business_analysis.tex, pitch-deck-demium.tex) and unchanged since: a
% document from any repository should look like it came from the same shop.
% ----------------------------------------------------------------------------
\definecolor{primaryblue}{RGB}{41,128,185}
\definecolor{darkblue}{RGB}{31,78,121}
\definecolor{lightgray}{RGB}{236,240,241}
\definecolor{darkgray}{RGB}{52,73,94}
\definecolor{accentgreen}{RGB}{39,174,96}
\definecolor{accentorange}{RGB}{230,126,34}
\definecolor{accentred}{RGB}{231,76,60}

% ----------------------------------------------------------------------------
% Section formatting.
% ----------------------------------------------------------------------------
\titleformat{\section}
    {\color{primaryblue}\Large\bfseries\sffamily}
    {\thesection.}{0.5em}{}[\titlerule]

\titleformat{\subsection}
    {\color{darkblue}\large\bfseries\sffamily}
    {\thesubsection}{0.5em}{}

\titleformat{\subsubsection}
    {\color{darkgray}\normalsize\bfseries\sffamily}
    {\thesubsubsection}{0.5em}{}

% ----------------------------------------------------------------------------
% Header and footer. \paperstatus rides in the right-hand header of every
% page, so a draft cannot be mistaken for a final document at any page a
% reader happens to open it at -- which is the whole reason the marker is in
% the header rather than only on the title page.
% ----------------------------------------------------------------------------
\pagestyle{fancy}
\fancyhf{}
\renewcommand{\headrulewidth}{0.4pt}
\fancyhead[L]{\small\sffamily\color{darkgray} \headerleft}
\fancyhead[R]{\small\sffamily\color{accentred} \paperstatus\ --- \today}
\fancyfoot[C]{\thepage}

% ----------------------------------------------------------------------------
% Links and PDF metadata.
% ----------------------------------------------------------------------------
\hypersetup{
    colorlinks=true,
    linkcolor=primaryblue,
    urlcolor=primaryblue,
    pdfauthor={\pdfauthorline},
    pdftitle={\pdftitleline}
}

% ----------------------------------------------------------------------------
% Figures. A diagram's source is one .mmd file rendered to vector PDF by the
% repository's render step; the document includes the rendered file.
%
% \dgm resolves the path from a slug and the edition, so the body text names a
% diagram the way the diagrams directory does and never writes a path:
%
%   \includegraphics[width=0.78\linewidth]{\dgm{a1-system-context}}
%
% resolves to ../diagrams/rendered/a1-system-context.pdf in the default
% edition and ../diagrams/rendered/a1-system-context.pl.pdf in an edition that
% set \def\editionsuffix{.pl}. A repository whose editions share one set of
% diagrams simply never sets \editionsuffix, and every edition resolves to the
% same file.
%
% Hardcoding the path instead is what the estate did first, and it costs one
% edit per figure per edition: `ab-ove` maintained `b1-reader-loop.pdf` in the
% English file against `b1-reader-loop.pl.pdf` in the Polish one, by hand, in
% both. At three figures that is survivable; `agent-eval-bench` has 23.
%
% \housediagramdir is overridable for a document that does not sit one level
% below the diagrams directory -- define it before the \input.
% ----------------------------------------------------------------------------
\providecommand{\editionsuffix}{}
\providecommand{\housediagramdir}{../diagrams/rendered}
\newcommand{\dgm}[1]{\housediagramdir/#1\editionsuffix.pdf}

% ----------------------------------------------------------------------------
% House convenience commands.
%
% \repopath is for a path, a filename or an identifier copied from the code;
% it is deliberately the only one of these, so that a document does not grow a
% private synonym for it. \finding marks the sentence a section exists to
% deliver.
% ----------------------------------------------------------------------------
\newcommand{\repopath}[1]{\colorbox{lightgray}{\texttt{#1}}}
\newcommand{\finding}[1]{{\color{primaryblue}\textbf{#1}}}

% ----------------------------------------------------------------------------
% The house title block, as a command so that every document's first page
% looks the same and none of them reimplements it.
%
%   \housetitle{Title}{Subtitle or empty}{Author}{Affiliation or repo URL}{Provenance line}
%
% The provenance line is the one this estate keeps paying for the absence of:
% a document that presents another document names it here, so a reader who
% finds the two disagreeing knows which one is the source. For a study paper
% it names the companion study and the commit; for a project overview, the
% markdown it presents.
% ----------------------------------------------------------------------------
\newcommand{\housetitle}[5]{%
    \thispagestyle{plain}%
    \begin{center}
        {\LARGE\bfseries\sffamily\color{darkblue} #1}

        \ifx\relax#2\relax\else
            \vspace{0.3cm}

            {\Large\sffamily\color{darkgray} #2}
        \fi

        \vspace{0.6cm}

        {\large #3}

        \vspace{0.2cm}

        {\normalsize\color{darkgray} #4}

        \vspace{0.2cm}

        {\normalsize\color{accentred}\sffamily \paperstatus\ --- \today\ --- #5}
    \end{center}

    \vspace{0.4cm}
}

%   \begin{document}
%   ...
%
% Its Polish sibling differs by two lines, both before the \input:
%
%   \def\houselang{polish}
%   \def\editionsuffix{.pl}
%
% ----------------------------------------------------------------------------
% PACKAGES THIS FILE DOES NOT LOAD. The house set below is what every document
% needs. A document that needs more -- tikz, pifont, longtable, fancyvrb,
% amssymb -- loads it in its own preamble AFTER the \input, and says in a
% comment that it is an addition rather than part of the house set, so the
% next reader can tell the two apart at a glance.
% ============================================================================

% ----------------------------------------------------------------------------
% Encoding, then language. fontenc before inputenc before babel: babel reads
% the active font encoding when it sets up its shorthands, so loading it first
% gives an edition that is translated but typeset wrong.
% ----------------------------------------------------------------------------
\usepackage[T1]{fontenc}
\usepackage[utf8]{inputenc}

% The edition's language, defaulted so a monolingual document defines nothing.
% \PassOptionsToPackage rather than \usepackage[\houselang]{babel}: the option
% list here is a macro, and PassOptionsToPackage is the mechanism built for
% passing one.
\providecommand{\houselang}{english}
\PassOptionsToPackage{\houselang}{babel}
\usepackage{babel}

% ----------------------------------------------------------------------------
% The house package set.
% ----------------------------------------------------------------------------
\usepackage{geometry}
\usepackage{xcolor}
\usepackage{hyperref}
\usepackage{graphicx}
\usepackage{enumitem}
\usepackage{fancyhdr}
\usepackage{titlesec}
\usepackage{parskip}
\usepackage{booktabs}
\usepackage{array}
\usepackage{tabularx}

% ----------------------------------------------------------------------------
% Page geometry.
% ----------------------------------------------------------------------------
\geometry{
    a4paper,
    left=2.2cm,
    right=2.2cm,
    top=2.5cm,
    bottom=2.5cm
}

% ----------------------------------------------------------------------------
% House colors. Extracted from the estate's first formal documents
% (business_analysis.tex, pitch-deck-demium.tex) and unchanged since: a
% document from any repository should look like it came from the same shop.
% ----------------------------------------------------------------------------
\definecolor{primaryblue}{RGB}{41,128,185}
\definecolor{darkblue}{RGB}{31,78,121}
\definecolor{lightgray}{RGB}{236,240,241}
\definecolor{darkgray}{RGB}{52,73,94}
\definecolor{accentgreen}{RGB}{39,174,96}
\definecolor{accentorange}{RGB}{230,126,34}
\definecolor{accentred}{RGB}{231,76,60}

% ----------------------------------------------------------------------------
% Section formatting.
% ----------------------------------------------------------------------------
\titleformat{\section}
    {\color{primaryblue}\Large\bfseries\sffamily}
    {\thesection.}{0.5em}{}[\titlerule]

\titleformat{\subsection}
    {\color{darkblue}\large\bfseries\sffamily}
    {\thesubsection}{0.5em}{}

\titleformat{\subsubsection}
    {\color{darkgray}\normalsize\bfseries\sffamily}
    {\thesubsubsection}{0.5em}{}

% ----------------------------------------------------------------------------
% Header and footer. \paperstatus rides in the right-hand header of every
% page, so a draft cannot be mistaken for a final document at any page a
% reader happens to open it at -- which is the whole reason the marker is in
% the header rather than only on the title page.
% ----------------------------------------------------------------------------
\pagestyle{fancy}
\fancyhf{}
\renewcommand{\headrulewidth}{0.4pt}
\fancyhead[L]{\small\sffamily\color{darkgray} \headerleft}
\fancyhead[R]{\small\sffamily\color{accentred} \paperstatus\ --- \today}
\fancyfoot[C]{\thepage}

% ----------------------------------------------------------------------------
% Links and PDF metadata.
% ----------------------------------------------------------------------------
\hypersetup{
    colorlinks=true,
    linkcolor=primaryblue,
    urlcolor=primaryblue,
    pdfauthor={\pdfauthorline},
    pdftitle={\pdftitleline}
}

% ----------------------------------------------------------------------------
% Figures. A diagram's source is one .mmd file rendered to vector PDF by the
% repository's render step; the document includes the rendered file.
%
% \dgm resolves the path from a slug and the edition, so the body text names a
% diagram the way the diagrams directory does and never writes a path:
%
%   \includegraphics[width=0.78\linewidth]{\dgm{a1-system-context}}
%
% resolves to ../diagrams/rendered/a1-system-context.pdf in the default
% edition and ../diagrams/rendered/a1-system-context.pl.pdf in an edition that
% set \def\editionsuffix{.pl}. A repository whose editions share one set of
% diagrams simply never sets \editionsuffix, and every edition resolves to the
% same file.
%
% Hardcoding the path instead is what the estate did first, and it costs one
% edit per figure per edition: `ab-ove` maintained `b1-reader-loop.pdf` in the
% English file against `b1-reader-loop.pl.pdf` in the Polish one, by hand, in
% both. At three figures that is survivable; `agent-eval-bench` has 23.
%
% \housediagramdir is overridable for a document that does not sit one level
% below the diagrams directory -- define it before the \input.
% ----------------------------------------------------------------------------
\providecommand{\editionsuffix}{}
\providecommand{\housediagramdir}{../diagrams/rendered}
\newcommand{\dgm}[1]{\housediagramdir/#1\editionsuffix.pdf}

% ----------------------------------------------------------------------------
% House convenience commands.
%
% \repopath is for a path, a filename or an identifier copied from the code;
% it is deliberately the only one of these, so that a document does not grow a
% private synonym for it. \finding marks the sentence a section exists to
% deliver.
% ----------------------------------------------------------------------------
\newcommand{\repopath}[1]{\colorbox{lightgray}{\texttt{#1}}}
\newcommand{\finding}[1]{{\color{primaryblue}\textbf{#1}}}

% ----------------------------------------------------------------------------
% The house title block, as a command so that every document's first page
% looks the same and none of them reimplements it.
%
%   \housetitle{Title}{Subtitle or empty}{Author}{Affiliation or repo URL}{Provenance line}
%
% The provenance line is the one this estate keeps paying for the absence of:
% a document that presents another document names it here, so a reader who
% finds the two disagreeing knows which one is the source. For a study paper
% it names the companion study and the commit; for a project overview, the
% markdown it presents.
% ----------------------------------------------------------------------------
\newcommand{\housetitle}[5]{%
    \thispagestyle{plain}%
    \begin{center}
        {\LARGE\bfseries\sffamily\color{darkblue} #1}

        \ifx\relax#2\relax\else
            \vspace{0.3cm}

            {\Large\sffamily\color{darkgray} #2}
        \fi

        \vspace{0.6cm}

        {\large #3}

        \vspace{0.2cm}

        {\normalsize\color{darkgray} #4}

        \vspace{0.2cm}

        {\normalsize\color{accentred}\sffamily \paperstatus\ --- \today\ --- #5}
    \end{center}

    \vspace{0.4cm}
}

%   \begin{document}
%   ...
%
% Its Polish sibling differs by two lines, both before the \input:
%
%   \def\houselang{polish}
%   \def\editionsuffix{.pl}
%
% ----------------------------------------------------------------------------
% PACKAGES THIS FILE DOES NOT LOAD. The house set below is what every document
% needs. A document that needs more -- tikz, pifont, longtable, fancyvrb,
% amssymb -- loads it in its own preamble AFTER the \input, and says in a
% comment that it is an addition rather than part of the house set, so the
% next reader can tell the two apart at a glance.
% ============================================================================

% ----------------------------------------------------------------------------
% Encoding, then language. fontenc before inputenc before babel: babel reads
% the active font encoding when it sets up its shorthands, so loading it first
% gives an edition that is translated but typeset wrong.
% ----------------------------------------------------------------------------
\usepackage[T1]{fontenc}
\usepackage[utf8]{inputenc}

% The edition's language, defaulted so a monolingual document defines nothing.
% \PassOptionsToPackage rather than \usepackage[\houselang]{babel}: the option
% list here is a macro, and PassOptionsToPackage is the mechanism built for
% passing one.
\providecommand{\houselang}{english}
\PassOptionsToPackage{\houselang}{babel}
\usepackage{babel}

% ----------------------------------------------------------------------------
% The house package set.
% ----------------------------------------------------------------------------
\usepackage{geometry}
\usepackage{xcolor}
\usepackage{hyperref}
\usepackage{graphicx}
\usepackage{enumitem}
\usepackage{fancyhdr}
\usepackage{titlesec}
\usepackage{parskip}
\usepackage{booktabs}
\usepackage{array}
\usepackage{tabularx}

% ----------------------------------------------------------------------------
% Page geometry.
% ----------------------------------------------------------------------------
\geometry{
    a4paper,
    left=2.2cm,
    right=2.2cm,
    top=2.5cm,
    bottom=2.5cm
}

% ----------------------------------------------------------------------------
% House colors. Extracted from the estate's first formal documents
% (business_analysis.tex, pitch-deck-demium.tex) and unchanged since: a
% document from any repository should look like it came from the same shop.
% ----------------------------------------------------------------------------
\definecolor{primaryblue}{RGB}{41,128,185}
\definecolor{darkblue}{RGB}{31,78,121}
\definecolor{lightgray}{RGB}{236,240,241}
\definecolor{darkgray}{RGB}{52,73,94}
\definecolor{accentgreen}{RGB}{39,174,96}
\definecolor{accentorange}{RGB}{230,126,34}
\definecolor{accentred}{RGB}{231,76,60}

% ----------------------------------------------------------------------------
% Section formatting.
% ----------------------------------------------------------------------------
\titleformat{\section}
    {\color{primaryblue}\Large\bfseries\sffamily}
    {\thesection.}{0.5em}{}[\titlerule]

\titleformat{\subsection}
    {\color{darkblue}\large\bfseries\sffamily}
    {\thesubsection}{0.5em}{}

\titleformat{\subsubsection}
    {\color{darkgray}\normalsize\bfseries\sffamily}
    {\thesubsubsection}{0.5em}{}

% ----------------------------------------------------------------------------
% Header and footer. \paperstatus rides in the right-hand header of every
% page, so a draft cannot be mistaken for a final document at any page a
% reader happens to open it at -- which is the whole reason the marker is in
% the header rather than only on the title page.
% ----------------------------------------------------------------------------
\pagestyle{fancy}
\fancyhf{}
\renewcommand{\headrulewidth}{0.4pt}
\fancyhead[L]{\small\sffamily\color{darkgray} \headerleft}
\fancyhead[R]{\small\sffamily\color{accentred} \paperstatus\ --- \today}
\fancyfoot[C]{\thepage}

% ----------------------------------------------------------------------------
% Links and PDF metadata.
% ----------------------------------------------------------------------------
\hypersetup{
    colorlinks=true,
    linkcolor=primaryblue,
    urlcolor=primaryblue,
    pdfauthor={\pdfauthorline},
    pdftitle={\pdftitleline}
}

% ----------------------------------------------------------------------------
% Figures. A diagram's source is one .mmd file rendered to vector PDF by the
% repository's render step; the document includes the rendered file.
%
% \dgm resolves the path from a slug and the edition, so the body text names a
% diagram the way the diagrams directory does and never writes a path:
%
%   \includegraphics[width=0.78\linewidth]{\dgm{a1-system-context}}
%
% resolves to ../diagrams/rendered/a1-system-context.pdf in the default
% edition and ../diagrams/rendered/a1-system-context.pl.pdf in an edition that
% set \def\editionsuffix{.pl}. A repository whose editions share one set of
% diagrams simply never sets \editionsuffix, and every edition resolves to the
% same file.
%
% Hardcoding the path instead is what the estate did first, and it costs one
% edit per figure per edition: `ab-ove` maintained `b1-reader-loop.pdf` in the
% English file against `b1-reader-loop.pl.pdf` in the Polish one, by hand, in
% both. At three figures that is survivable; `agent-eval-bench` has 23.
%
% \housediagramdir is overridable for a document that does not sit one level
% below the diagrams directory -- define it before the \input.
% ----------------------------------------------------------------------------
\providecommand{\editionsuffix}{}
\providecommand{\housediagramdir}{../diagrams/rendered}
\newcommand{\dgm}[1]{\housediagramdir/#1\editionsuffix.pdf}

% ----------------------------------------------------------------------------
% House convenience commands.
%
% \repopath is for a path, a filename or an identifier copied from the code;
% it is deliberately the only one of these, so that a document does not grow a
% private synonym for it. \finding marks the sentence a section exists to
% deliver.
% ----------------------------------------------------------------------------
\newcommand{\repopath}[1]{\colorbox{lightgray}{\texttt{#1}}}
\newcommand{\finding}[1]{{\color{primaryblue}\textbf{#1}}}

% ----------------------------------------------------------------------------
% The house title block, as a command so that every document's first page
% looks the same and none of them reimplements it.
%
%   \housetitle{Title}{Subtitle or empty}{Author}{Affiliation or repo URL}{Provenance line}
%
% The provenance line is the one this estate keeps paying for the absence of:
% a document that presents another document names it here, so a reader who
% finds the two disagreeing knows which one is the source. For a study paper
% it names the companion study and the commit; for a project overview, the
% markdown it presents.
% ----------------------------------------------------------------------------
\newcommand{\housetitle}[5]{%
    \thispagestyle{plain}%
    \begin{center}
        {\LARGE\bfseries\sffamily\color{darkblue} #1}

        \ifx\relax#2\relax\else
            \vspace{0.3cm}

            {\Large\sffamily\color{darkgray} #2}
        \fi

        \vspace{0.6cm}

        {\large #3}

        \vspace{0.2cm}

        {\normalsize\color{darkgray} #4}

        \vspace{0.2cm}

        {\normalsize\color{accentred}\sffamily \paperstatus\ --- \today\ --- #5}
    \end{center}

    \vspace{0.4cm}
}


% Added for this document (not in the house set): a long table that may break
% across pages, and \texttt line breaking at the characters identifiers use.
\usepackage{longtable}
\usepackage[htt]{hyphenat}
% Keep each figure inside the section that discusses it.
\usepackage[section]{placeins}
% Long identifiers in running text: allow looser lines rather than overfull ones.
\setlength{\emergencystretch}{3em}

\begin{document}

\housetitle{payments-bench}{An outbound transfer service that moves money exactly once}{Project overview}{\url{https://github.com/konradcinkusz/payments-bench}}{presents docs/OVERVIEW.md}

\begin{abstract}
\noindent
payments-bench is an outbound bank-transfer service built to show the mechanisms a payment rail
needs to move money exactly once and to recover when the bank on the other end misbehaves:
idempotent initiation, a six-state payment machine, a submit whose status is committed before the
partner call and whose failures become an Unknown outcome that is never retried blind, a resolver
that asks the partner before it acts, a transactional outbox and an inbox, a double-entry
append-only ledger, and reconciliation that returns named breaks as a work queue. The bank is
simulated; the mechanisms are real code with tests. 37 tests pass (\repopath{tests/}, counted by
\texttt{dotnet test} on 2026-10-13). Five defects are recorded in the findings report, none of
them caught by a unit test, and ten rows of the register say what has not been proved, the first
being that nothing here has ever talked to a real bank.
\end{abstract}

\tableofcontents
\newpage

% ----------------------------------------------------------------------------
\section{Executive summary}

\finding{This is not a payments product. It is a bench for the mechanisms one depends on.}
Moving money is easy to get almost right. A transfer request arrives, a row is written, a call
goes to the bank, the bank says yes. Every one of those steps can happen twice, or half-happen, or
happen with nobody hearing the answer, and in each case naive code either loses track of a
customer's money or sends it twice. The repository builds the smallest service that handles those
cases on purpose, and writes down in a register which of them it has proved and which it has not.

What is on \texttt{main}:
\begin{itemize}[leftmargin=1.4em]
  \item an HTTP API (\repopath{src/PaymentsBench.Api/Program.cs}) to initiate, submit, resolve and
        inspect a transfer, receive partner notifications, run reconciliation and read balances;
  \item a six-state payment machine whose allowed moves are one table,
        \repopath{PaymentTransitions};
  \item idempotent initiation by \texttt{Idempotency-Key} plus a SHA-256 hash of the body;
  \item a submit that commits \texttt{Submitted} before calling the partner and maps every
        failure after the request may have left the process to an \textbf{Unknown} outcome;
  \item a resolver that asks the partner by end-to-end id and resubmits only on
        \texttt{NotFound};
  \item a transactional outbox for events and an inbox that deduplicates notifications;
  \item a double-entry, append-only ledger over \texttt{customer-funds},
        \texttt{pending-outbound} and \texttt{nostro};
  \item reconciliation of a day's statement against settled payments, as named breaks;
  \item scheme rules as data in \repopath{scenarios/schemes.yaml};
  \item a \textbf{simulated} partner bank with five fault profiles switched at runtime;
  \item a degraded mode: without a database connection string every store is in memory and the
        log says so.
\end{itemize}

The 37 tests run against in-memory stores and a scripted partner, which is exactly why the two
most serious defects found so far were found by walking a transfer through the running services
over HTTP instead (\finding{F-1}, \finding{F-2}; Section~\ref{sec:caught}).

% ----------------------------------------------------------------------------
\section{Why a payment rail is hard}

Four properties make outbound payments harder than ordinary request handling, and each maps to a
mechanism in the repository.

\begin{description}[leftmargin=1.4em,style=nextline]
  \item[Requests are repeated.] A client retries on a dropped connection, a queue redelivers, a
        user double-clicks. Without a stable notion of ``the same request'' each repeat is a new
        transfer. $\rightarrow$ idempotent initiation (Section~\ref{sec:idem}).
  \item[The partner's answer can be lost after the partner acted.] From the caller's side, a
        recorded transfer whose response was lost looks identical to a request that never
        arrived. A retry is safe in one case and a duplicate in the other, and the caller cannot
        tell which. $\rightarrow$ the Unknown outcome and the resolver
        (Section~\ref{sec:unknown}).
  \item[Delivery between systems is at least once.] Events and notifications can both arrive
        twice. $\rightarrow$ outbox and inbox (Section~\ref{sec:outbox}).
  \item[The money must be explainable afterwards.] ``Where is this customer's money right now''
        must have one answer derived from history, and what the bank's statement says must be
        compared with what the service believes. $\rightarrow$ the ledger
        (Section~\ref{sec:ledger}) and reconciliation (Section~\ref{sec:recon}).
\end{description}

Underneath sits a duller property: the schemes that carry the money have cut-offs, settlement
days and calendars that change independently of the code. $\rightarrow$ schemes as data
(Section~\ref{sec:schemes}).

% ----------------------------------------------------------------------------
\section{The money path}

Every movement of money is a balanced journal entry between three accounts. The postings for each
payment event are in one file, \repopath{src/PaymentsBench.Domain/Ledger/Postings.cs}, so the path
can be read top to bottom.

\begin{center}
\begin{tabular}{@{}llll@{}}
\toprule
Event & Posting & Debit & Credit \\
\midrule
Initiated & \texttt{Reserve} & \texttt{customer-funds} & \texttt{pending-outbound} \\
Settled & \texttt{Settle} & \texttt{pending-outbound} & \texttt{nostro} \\
Rejected (before settlement) & \texttt{Release} & \texttt{pending-outbound} & \texttt{customer-funds} \\
Returned (after settlement) & \texttt{Reverse} & \texttt{nostro} & \texttt{customer-funds} \\
\bottomrule
\end{tabular}
\end{center}

\begin{figure}[!htbp]
  \centering
  \includegraphics[width=0.82\linewidth,height=0.5\textheight,keepaspectratio]{\dgm{a1-money-path}}
  \caption{The money path: which payment event posts which movement between the three ledger
  accounts (\repopath{docs/diagrams/a1-money-path.mmd}).}
\end{figure}

Customer funds are reserved the moment a payment is initiated, sit in \texttt{pending-outbound}
while the transfer is in flight, and leave through \texttt{nostro} only when the partner says the
transfer settled. A rejection releases the reservation; a return reverses the settlement.
\finding{After reserve, settle and reverse, the customer's money is back where it started and
every account nets to zero}; the ledger test
\texttt{The\_full\_money\_path\_of\_a\_returned\_payment\_nets\_to\_zero\_on\_customer\_funds\_and\_nostro}
proves it.

% ----------------------------------------------------------------------------
\section{Architecture}

Four projects carry the logic, in strict dependency order.

\begin{center}
\begin{tabularx}{\linewidth}{@{}l l X@{}}
\toprule
Project & Knows & Holds \\
\midrule
\texttt{Domain} & nothing & \texttt{Money}, \texttt{Payment}, \texttt{PaymentTransitions}, \texttt{JournalEntry}, \texttt{Postings}, \texttt{SchemeRules}, \texttt{ValueDateCalculator} \\
\texttt{Application} & domain, ports & the handlers (initiate, submit, resolve, apply notification), \texttt{ReconciliationService}, the ports \\
\texttt{Infrastructure} & the ports & in-memory and EF Core/PostgreSQL stores, \texttt{OutboxDispatcher}, \texttt{HttpPartnerGateway}, \texttt{HttpStatementSource} \\
\texttt{Api} & all of the above & the HTTP surface; \texttt{Program.cs} is the manifest \\
\bottomrule
\end{tabularx}
\end{center}

Around them, \texttt{ServiceDefaults} carries OpenTelemetry, health and service discovery and by
rule no domain type; \texttt{AppHost} is the Aspire composition root (one PostgreSQL with a
\texttt{payments} database, the partner simulator, the API); \texttt{PartnerSimulator} is the
simulated bank.

\begin{figure}[!htbp]
  \centering
  \includegraphics[width=0.82\linewidth,height=0.5\textheight,keepaspectratio]{\dgm{a6-delivery-topology}}
  \caption{What the composition root declares, and the degraded in-memory mode
  (\repopath{docs/diagrams/a6-delivery-topology.mmd}).}
\end{figure}

\paragraph{Storage modes.} \repopath{InfrastructureRegistration} picks the stores once, at
start-up. With a connection string named \texttt{payments}, every store is EF Core on PostgreSQL,
the outbox has a dispatcher, and the schema is applied by a hosted service. Without one, every
store is in memory, the unit of work is a no-op, and the log says at warning level that the
service is running degraded and nothing survives a restart. \finding{A missing database degrades
the service instead of stopping it}; the tests use this mode.

\paragraph{Observability.} One OpenTelemetry trace per request; a \texttt{payments.transitions}
counter tagged by status; \texttt{/health} and \texttt{/alive}; OTLP export when
\texttt{OTEL\_EXPORTER\_OTLP\_ENDPOINT} is set.

% ----------------------------------------------------------------------------
\section{The seven rules}

\begin{center}
\begin{tabularx}{\linewidth}{@{}r X X@{}}
\toprule
\# & Rule & Enforced in \\
\midrule
1 & Submitted is committed before the partner call. & \texttt{SubmitPaymentHandler}: mark, update, outbox, commit, then call \\
2 & An unknown outcome is never retried. & \texttt{HttpPartnerGateway} maps failures to \texttt{Unknown}; the pipeline retries \texttt{GET} only \\
3 & Resubmission needs the partner's word. & \texttt{ResolveSubmittedPaymentHandler}: only \texttt{NotFound}, same end-to-end id \\
4 & The ledger is append-only and balanced. & \texttt{JournalEntry} cannot be built unbalanced per currency \\
5 & Events through the outbox, notifications through the inbox. & every handler's unit of work; \texttt{IInbox.TryRecordAsync} first \\
6 & Reconciliation produces breaks, not alerts. & \texttt{ReconciliationReport} with named \texttt{BreakKind}s \\
7 & Scheme rules are data. & \repopath{scenarios/schemes.yaml} \\
\bottomrule
\end{tabularx}
\end{center}

Rules 1 to 3 are one argument and the core of the repository. \finding{If the state before the
call is durable, an unknowable result is never guessed at, and the only thing allowed to trigger
a second attempt is the partner's own statement that the first one does not exist, then a
transfer is sent at most once and is never lost track of.} \finding{F-2} is the evidence that
the argument holds even when the code around it is wrong.

% ----------------------------------------------------------------------------
\section{Idempotency}\label{sec:idem}

\texttt{POST /payments} requires an \texttt{Idempotency-Key} header; without it the API answers
400. \texttt{InitiatePaymentHandler} computes a SHA-256 over a canonical string of the fields that
make two requests the same request: end-to-end id, amount to two decimals in the invariant
culture, currency upper-cased, debtor and creditor accounts, and scheme. No payment with that key:
create one, reserve the funds and enqueue \texttt{PaymentInitiated} in one unit of work. Same key,
same hash: return the original. Same key, different hash: a 409. Three tests pin this. Two
identical requests racing into the database at once are not yet tested; that sits with register
row R6.

% ----------------------------------------------------------------------------
\section{The unknown outcome and the resolver}\label{sec:unknown}

ADR 0001 is the decision the repository is built around. A submit can fail in three ways that
look alike: the connection drops before the request lands, the request lands and the response is
lost, or the partner records the transfer and returns a 5xx. \finding{A naive retry creates a
second transfer in two of the three.}

\begin{figure}[!htbp]
  \centering
  \includegraphics[width=0.95\linewidth,height=0.5\textheight,keepaspectratio]{\dgm{a3-submit-resolve}}
  \caption{Submit and resolve: \texttt{Submitted} committed before the call; Accepted, Rejected
  and Unknown; the resolver asks by end-to-end id and resubmits only on \texttt{NotFound}
  (\repopath{docs/diagrams/a3-submit-resolve.mmd}).}
\end{figure}

\paragraph{Submit.} The handler moves the payment to \texttt{Submitted}, enqueues the event and
commits, and only then calls the partner. The gateway maps a 422 to \texttt{Rejected}; any other
non-success status, an empty body, or an \texttt{HttpRequestException},
\texttt{TaskCanceledException}, \texttt{TimeoutRejectedException} or
\texttt{BrokenCircuitException} to \texttt{Unknown}; a success with a partner reference to
\texttt{Accepted}. \finding{On Unknown the handler does nothing at all}: the payment stays at
\texttt{Submitted}, the state that means ``ask before you act''. The resilience handler is
configured to retry only \texttt{GET}, so the submit is sent once; its predicate is stricter than
intended and never retries a \texttt{GET} on an exception either (\finding{F-5}).

\paragraph{Resolve.} If the payment is not at \texttt{Submitted}, nothing happens. Otherwise the
resolver asks the partner by end-to-end id: \texttt{NotFound} $\rightarrow$ resubmit once with
the same end-to-end id; \texttt{Accepted}, \texttt{Settled}, \texttt{Rejected} $\rightarrow$ apply
it with its posting; \texttt{Returned} $\rightarrow$ refuse, since a return for a transfer never
accepted on this side is an error. The same end-to-end id lets the partner recognise a late first
attempt as a duplicate; that is an assumption about the partner, stated in ADR 0001, and the simulated partner honours it by design. Five tests cover the resolver's decisions. One branch is untested
and named badly: when the resubmit's own outcome is unknown, the resolver returns
\texttt{SafeToResubmit} (\finding{F-3}). Nothing calls the resolver on its own yet (R8).

\begin{figure}[!htbp]
  \centering
  \includegraphics[width=0.62\linewidth,height=0.5\textheight,keepaspectratio]{\dgm{a2-state-machine}}
  \caption{The six states and every allowed transition, as listed in
  \repopath{PaymentTransitions} (\repopath{docs/diagrams/a2-state-machine.mmd}).}
\end{figure}

% ----------------------------------------------------------------------------
\section{Outbox and inbox}\label{sec:outbox}

Every handler enqueues the payment's domain events through \texttt{IOutbox} in the same unit of
work as the state change and the ledger entry. With PostgreSQL, \texttt{OutboxDispatcher} polls
every two seconds, takes up to 100 undispatched messages oldest first, publishes each and marks it
dispatched only after the publish returned: \finding{delivery is at least once by design}, and
consumers must deduplicate on the message id. It publishes to a log line today; no broker and no
consumer exist (R5). In the in-memory mode there is no dispatcher at all.

Partner notifications are recorded in \texttt{IInbox} by notification id before anything else; a
second delivery returns \texttt{Duplicate} and changes nothing. A notification for an unknown
end-to-end id is recorded and reported as \texttt{UnknownPayment}, not dropped, because it is what
reconciliation needs to see.

\begin{figure}[!htbp]
  \centering
  \includegraphics[width=0.85\linewidth,height=0.5\textheight,keepaspectratio]{\dgm{a4-outbox-inbox}}
  \caption{Outbox and inbox (\repopath{docs/diagrams/a4-outbox-inbox.mmd}).}
\end{figure}

% ----------------------------------------------------------------------------
\section{The ledger}\label{sec:ledger}

ADR 0002: every money movement is a \texttt{JournalEntry} of at least two lines that balance per
currency; entries are appended, never updated or deleted; a correction is another entry; balances
are sums over lines; direction is the side, never a negative amount. Five ledger tests prove that
an unbalanced entry cannot exist, that balance is per currency, that direction is the side, that a
rejection releases the reservation, and that a returned payment nets to zero. Balances are sums
over rows; at volume that needs a materialised balance maintained by the same entries, which stay
the source of truth.

% ----------------------------------------------------------------------------
\section{Reconciliation}\label{sec:recon}

ADR 0003: a statement that disagrees with the ledger is the normal case on some days, and an
alert saying ``reconciliation failed'' tells the person on call nothing. The service matches the
day's debit lines to payments that are \texttt{Settled} with that value date, by end-to-end id and
exact amount and currency, and returns the number matched and a list of breaks, each with both
amounts where both exist.

\begin{center}
\begin{tabularx}{\linewidth}{@{}l X@{}}
\toprule
Break kind & Meaning \\
\midrule
\texttt{OnStatementOnly} & the statement books a movement with no settled payment behind it \\
\texttt{InLedgerOnly} & a payment settled that the statement does not show \\
\texttt{AmountMismatch} & both sides have it, with different amounts \\
\bottomrule
\end{tabularx}
\end{center}

\begin{figure}[!htbp]
  \centering
  \includegraphics[width=0.85\linewidth,height=0.5\textheight,keepaspectratio]{\dgm{a5-reconciliation}}
  \caption{Reconciliation as a work queue (\repopath{docs/diagrams/a5-reconciliation.mmd}).}
\end{figure}

One defect is known: a payment settled on a day and returned later is no longer
\texttt{Settled}, so its settlement line on that day's statement is reported as
\texttt{OnStatementOnly} (\finding{F-4}). The simulated statement is a JSON reduction of what an
ISO 20022 \texttt{camt.053} carries and has no line for the return itself (R2).

% ----------------------------------------------------------------------------
\section{Schemes and calendars}\label{sec:schemes}

\begin{center}
\begin{tabular}{@{}llll@{}}
\toprule
Scheme & Currency & Cut-off & Settlement business days \\
\midrule
\texttt{sepa} & EUR & 16:00 & 1 \\
\texttt{faster-payments} & GBP & 23:59 & 0 \\
\texttt{swift} & USD & 15:00 & 2 \\
\bottomrule
\end{tabular}
\end{center}

The value date is the same business day before cut-off, else the next business day, plus the
scheme's settlement days (\repopath{scenarios/schemes.yaml}, \texttt{ValueDateCalculator}). The only
calendar is weekends-only: no public holidays, no TARGET2 closing days (R7).

% ----------------------------------------------------------------------------
\section{The partner simulator and the scenario classes}

\finding{The partner is a simulated bank whose API shape is invented} (R1). It accepts transfers
(deduplicated by end-to-end id), answers lookups, lets a test settle or return a transfer, serves
a statement per date, and switches its fault profile at runtime: \texttt{None},
\texttt{Latency} (1,500\,ms by default), \texttt{TimeoutAfterAccept} (records the transfer, holds
the response 10,000\,ms by default), \texttt{ServerError} (503, records nothing),
\texttt{RejectAll} (422, \texttt{AC01} by default), \texttt{StatementDrift} (0.01 by default added
to every statement amount).

Seven scenario files in four classes describe what the rail must survive, as data: \texttt{happy}
(2), \texttt{partner-failure} (2), \texttt{duplicate} (2), \texttt{reconciliation} (1). \finding{They
are a specification, not an executed suite}: a test checks that they parse; no runner drives them
against the API and the simulator yet (R3). That runner is also the closing condition that would
guard F-1 and F-2 against regression.

% ----------------------------------------------------------------------------
\section{CI and governance}

\repopath{.github/workflows/ci.yml} builds and tests on every pull request and on pushes to
\texttt{main}, and lints the documentation. \repopath{AGENTS.md} sets the rules: .NET 10, warnings
as errors, one pull request per change with CI green, scheme rules and scenarios as data, nothing
that talks to a real bank, no secrets in the tree. The register is never silently wrong; a new
decision is a new record. Further work is planned in three streams under \repopath{docs/plans/},
and until a slice merges this document does not describe it as done.

% ----------------------------------------------------------------------------
\section{What the tests and the walk-through caught}\label{sec:caught}

\finding{Five findings, none caught by a unit test} (\repopath{docs/FINDINGS.md}).

\begin{center}
\begin{tabularx}{\linewidth}{@{}l X l l@{}}
\toprule
\# & Defect & Found by & Status \\
\midrule
F-1 & Enums crossed the wire as integers; every partner notification and statement line failed to bind & walk-through & fixed \\
F-2 & A timeout after the partner accepted surfaced as a 500, not as an Unknown outcome & walk-through & fixed \\
F-3 & The resolver reports \texttt{SafeToResubmit} after a resubmit whose own outcome is unknown & writing the overview & open \\
F-4 & A returned payment is an \texttt{OnStatementOnly} break for its settlement day & writing the overview & open \\
F-5 & The status lookup is never retried on an exception, only on a 5xx & writing the overview & open \\
\bottomrule
\end{tabularx}
\end{center}

The pattern is the finding. The tests are good at what they test, in one process, against
in-memory stores and a scripted partner; a defect between processes (F-1), in a library's
exception types (F-2), in an HTTP pipeline no test drives (F-5), or in an assumption the code and
its tests share (F-3, F-4) cannot fail them. \finding{F-2 is also the best evidence for the
design: the defect was real, and committing the status before the call meant it cost an error
message, not a duplicate transfer.}

% ----------------------------------------------------------------------------
\section{Production readiness, honestly}

\begin{itemize}[leftmargin=1.4em]
  \item \textbf{The partner is simulated.} No real bank, no scheme sandbox, no ISO 20022 message
        (R1, R2); lookup by end-to-end id is an assumption the simulator satisfies (ADR 0001).
  \item \textbf{The scenarios do not run} (R3); the two walk-through findings have no automated guard.
  \item \textbf{PostgreSQL has not run in CI}; the schema is applied with \texttt{EnsureCreated}
        and no migration exists (R4).
  \item \textbf{The degraded mode is degraded}: a no-op unit of work, no dispatcher, nothing
        survives a restart.
  \item \textbf{Concurrency is a column} (R6). \textbf{Nothing resolves on its own} (R8).
        \textbf{No load, no chaos} (R9).
  \item \textbf{No authentication on the API; no sanctions, AML or limits} (R10 for the controls).
  \item \textbf{Three open defects} (F-3, F-4, F-5).
\end{itemize}

% ----------------------------------------------------------------------------
\section{The register}

Ten open rows, none closed (\repopath{docs/REGISTER.md}), quoted by headline:

\begin{longtable}{@{}l p{0.82\linewidth}@{}}
\toprule
\# & Not proved \\
\midrule
\endhead
R1 & Nothing here has ever talked to a real bank. \\
R2 & No ISO 20022 message is parsed or produced. \\
R3 & The scenario files are a specification, not an executed suite. \\
R4 & PostgreSQL persistence has not run in CI. \\
R5 & The outbox dispatcher publishes to a log line. \\
R6 & Concurrency is a column, not a test. \\
R7 & The business calendar is weekends only. \\
R8 & The resolver is called, never scheduled. \\
R9 & No load, no chaos. \\
R10 & Sanctions, AML, limits: absent. \\
\bottomrule
\end{longtable}

% ----------------------------------------------------------------------------
\section{Decisions and non-goals}

\textbf{ADR 0001}: an unknown submit outcome is never retried blind. \textbf{ADR 0002}: money is a
double-entry, append-only ledger. \textbf{ADR 0003}: reconciliation produces breaks, not alerts.

Non-goals: being a payments product; talking to a real bank; compliance controls (R10); inbound
payments, FX, fees and netting; fuzzy reconciliation (anything else is a new, named break kind).

% ----------------------------------------------------------------------------
\section{Relationship to the standards}

The repository conforms to architecture-standards \cite{standards} and does not re-derive it. Its
documentation follows agent-eval-bench \cite{aeb}: a long-form overview as the source of truth, a
findings report written to be unflattering, a register of what is not proved, and presentation
editions, like this one, that add nothing the overview does not contain.

% ----------------------------------------------------------------------------
\section{Reproducibility}

Every number in this document is reproducible from the repository at the commit this PDF was
built from.

\begin{verbatim}
git clone https://github.com/konradcinkusz/payments-bench
cd payments-bench
dotnet test                 # the test counts: 15 domain + 22 application
npm ci                      # pinned doc tooling, including mermaid-cli
npm run lint                # markdown, links, diagrams, papers
npm run render:diagrams     # docs/diagrams/*.mmd -> docs/diagrams/rendered/*.pdf
cd docs/papers && latexmk -pdf payments-bench-overview.tex
\end{verbatim}

The walk-through behind F-1 and F-2 is the README's ``Run it'' section: start
\texttt{PaymentsBench.PartnerSimulator} and \texttt{PaymentsBench.Api} without a database, then
initiate, submit under \texttt{TimeoutAfterAccept}, resolve, settle, notify twice and reconcile with
\texttt{curl}. The F-4 probe is described step by step in the findings report.

% ----------------------------------------------------------------------------
\begin{thebibliography}{9}
\bibitem{overview} payments-bench, \texttt{docs/OVERVIEW.md}: the source of truth this document presents. \url{https://github.com/konradcinkusz/payments-bench/blob/main/docs/OVERVIEW.md}
\bibitem{findings} payments-bench, \texttt{docs/FINDINGS.md}: what building and testing the bench caught. \url{https://github.com/konradcinkusz/payments-bench/blob/main/docs/FINDINGS.md}
\bibitem{register} payments-bench, \texttt{docs/REGISTER.md}: what is not proved. \url{https://github.com/konradcinkusz/payments-bench/blob/main/docs/REGISTER.md}
\bibitem{arch} payments-bench, \texttt{docs/ARCHITECTURE.md}: the shape and the seven rules. \url{https://github.com/konradcinkusz/payments-bench/blob/main/docs/ARCHITECTURE.md}
\bibitem{adr} payments-bench, \texttt{docs/adr/}: decision records 0001--0003. \url{https://github.com/konradcinkusz/payments-bench/tree/main/docs/adr}
\bibitem{scenarios} payments-bench, \texttt{scenarios/}: the fault scenarios and scheme rules as data. \url{https://github.com/konradcinkusz/payments-bench/tree/main/scenarios}
\bibitem{standards} architecture-standards. \url{https://github.com/konradcinkusz/architecture-standards}
\bibitem{aeb} agent-eval-bench, the documentation model. \url{https://github.com/konradcinkusz/agent-eval-bench}
\end{thebibliography}

\end{document}
